% Revize 2026-10-28: obsah, matematika a návaznost; podrobnosti a časy
% viz planovani/REVIZE_PRIPRAV_1_ROCNIK.md.
% Aktuální porovnání se skutečným obsahem místních PDF Realisticky.cz:
% externi_zdroje/realisticky/01_Zakladni_poznatky/04_Formalni_logika_vyroky/010410_Formalni_logika_shrnuti/010410_Lekce.pdf
% externi_zdroje/realisticky/01_Zakladni_poznatky/02_Ciselne_obory/010205_Realna_cisla_I/010205_Lekce.pdf
% externi_zdroje/realisticky/01_Zakladni_poznatky/02_Ciselne_obory/010210_Absolutni_hodnota_I/010210_Lekce.pdf
% Rozpis: 1_rocnik/tematicke_plany/Rozpis_vyuky_MAT_1_rocnik.tex, hodina 10, T1.
% Rozsah: logika a číselné obory z hodin 02–09, nikoli množiny z 12 a dalších hodin.
% Učebnice: ucebnice/01_Zakladni_poznatky/Vyrokova_logika/Výroková logika.tex,
% Implikace, Negace složených výroků, Kvantifikované výroky a jejich negace;
% Ciselne_obory/Číselné obory.tex, Porovnávání racionálních čísel a Intervaly.
% Realisticky.cz: externi_zdroje/realisticky/01_Zakladni_poznatky/
% 04_Formalni_logika_vyroky/010410_Formalni_logika_shrnuti/010410_Lekce.pdf,
% úvodní shrnutí a Zádrhele; dále 02_Ciselne_obory/010203_Racionalni_cisla_I/
% 010203_Lekce.pdf (dělení nulou, zápisy). Modelové opravné úlohy jsou vlastní tvorba.
% Toto je pouze učitelská příprava, nikoli test ani návrh jeho konkrétní skladby.
\documentclass[11pt,a4paper]{article}
% Spolecny styl pro ucitelske pripravy MAT (Work/MAT)
% Verze 1.1, 2026-10-23
\NeedsTeXFormat{LaTeX2e}
\ProvidesPackage{priprava}[2026/10/23 v1.1 Sablona ucitelskych priprav MAT]

% Zakladni typografie a cestina
\RequirePackage[T1]{fontenc}
\RequirePackage[utf8]{inputenc}
\RequirePackage[czech]{babel}
\RequirePackage{lmodern}
\RequirePackage{microtype}
\RequirePackage{amsmath,amssymb}
\RequirePackage{xcolor}
\RequirePackage{enumitem}
\RequirePackage[most]{tcolorbox}
\RequirePackage{tabularx,booktabs,array}
\RequirePackage{fancyhdr}
\RequirePackage{needspace}
\RequirePackage{geometry}
\RequirePackage[hidelinks]{hyperref}

\geometry{a4paper,top=17mm,bottom=18mm,left=23mm,right=23mm}
\definecolor{notegray}{gray}{0.39}
\setlength{\parindent}{0pt}
\setlength{\parskip}{2.5pt plus 1pt minus 0.5pt}
\setlist[itemize]{topsep=2pt,itemsep=1pt,parsep=0pt,leftmargin=1.45em}
\setlist[enumerate]{topsep=2pt,itemsep=2pt,parsep=0pt,leftmargin=1.65em}
\setcounter{secnumdepth}{0}
\renewcommand{\arraystretch}{1.12}

% Zahlavi: vlevo rocnik | globalni tema, vpravo cislo: nazev hodiny
% Zapati zustava zcela prazdne.
\pagestyle{fancy}
\fancyhf{}
\renewcommand{\headrulewidth}{0.35pt}
\renewcommand{\footrulewidth}{0pt}
\setlength{\headheight}{14pt}
\newcommand{\nastavPripravu}[4]{%
  \fancyhead[L]{\small #1\enspace\textbar\enspace #2}%
  \fancyhead[R]{\small #3: #4}%
}

% Minimalisticke cernobile boxy. Kazda definice zvlast.
\tcbset{commonlesson/.style={enhanced,breakable=false,colback=white,colframe=black,
  colbacktitle=white,coltitle=black,boxrule=0.45pt,
  arc=0.8mm,outer arc=0.8mm,left=3mm,right=3mm,top=1.6mm,bottom=1.6mm,
  toptitle=1mm,bottomtitle=0.7mm,fonttitle=\bfseries,
  attach boxed title to top left={xshift=2.5mm,yshift=-1.15mm},
  boxed title style={colback=white,colframe=white,boxrule=0pt,left=1mm,right=1mm}}}
\newtcolorbox{zadaniBox}[1][]{commonlesson,title={Zadání},#1}
\newtcolorbox{prikladBox}[1][]{commonlesson,title={Příklad},#1}
\newtcolorbox{definiceBox}[1][]{commonlesson,title={Definice},#1}

% Spolecny box pro zadani a jeho reseni, bez mezery mezi castmi.
% Pouziti:
%   \begin{prikladReseniBox}
%     Zadani ulohy.
%     \tcblower
%     \begin{align*} ... \end{align*}
%   \end{prikladReseniBox}
% Spodni cast ma jemny sedy podklad a delici linku; lze v ni sazet
% vice radku matematiky. Obe casti maji jen jeden spolecny ramecek.
\tcbset{solutionlesson/.style={
  commonlesson,
  enhanced jigsaw,
  breakable,
  colbacklower=black!3,
  segmentation style={solid,draw=black!25,line width=0.35pt},
  fontlower=\small\color{notegray},
  before lower={\noindent\textbf{Řešení:}\enspace},
  leftlower=3mm,rightlower=3mm
}}
\newtcolorbox{prikladReseniBox}[1][]{solutionlesson,title={Příklad},#1}
\newtcolorbox{zadaniReseniBox}[1][]{solutionlesson,title={Zadání},#1}

% Pomucky pro autora pripravy: sedive poznamky; \reseni pro starsi pripravy
\newcommand{\poznamka}[1]{\par\begingroup\small\color{notegray}\itshape
  \textnormal{\textbf{Poznámka pro učitele.}}\enspace #1\par\endgroup}
\newcommand{\reseni}[1]{\par\begingroup\small\color{notegray}
  \textbf{Očekávané řešení:}\enspace #1\par\endgroup}
\newcommand{\blok}[1]{\needspace{5\baselineskip}\subsection*{#1}}
\newcommand{\vyrok}[1]{\uv{#1}}

\nastavPripravu{1. ročník}{Číselné obory}{10}{Průběžný test a rozbor}
\begin{document}

\blok{Organizace hodiny}
Před hodinou zajisti samostatně připravené a schválené zadání průběžného
testu z výrokové logiky a číselných oborů. Jeho vznik patří do postupu
\texttt{vytvor-test}; tato příprava žádná testová zadání neobsahuje.
Ověř soulad s hodinami 02–09: spojky, kvantifikátory a negace,
číselné obory, znaménka, zlomky, osa, zaokrouhlování, intervaly a význam
absolutní hodnoty. Zatím nepatří množinové operace ani obecné algoritmy
řešení rovnic s absolutní hodnotou.

\poznamka{Orientační průběh 45 min: organizace 2 min, test 18 min,
krátký rozbor 8 min, navazující procvičení 12 min, závěrečná kontrola
3 min a zadání domácí práce 2 min. Délku testu lze zvolit v rozmezí
15–20 min; následné procvičení tomu uprav. Rezervní blok níže je zásoba
na 10–15 min, kterou vybíráš místo části běžného procvičení, nikoli navíc
k celému programu. Domácí práce přibližně 10–15 min.}

Na začátku sděl délku testu, způsob odevzdání a pomůcky podle schváleného
zadání. Během testu poskytuj jen organizační vysvětlení, nikoli nápovědu
k řešení. Po vybrání zadání odděl hodnocenou práci od společné výuky.

\blok{Krátký rozbor: pojmenovat chybu a opravit postup}
Podle skutečně použitého zadání vyber dvě až tři důležité situace.
Neopisuj celé řešení testu; nech nejprve žáky říci, kde rozhodovalo
znaménko, obor proměnné nebo krajní bod. Pokud výsledky ještě nejsou
opravené, použij následující \emph{modelové} chyby a skutečná zjištění
z písemek přenes do hodiny 11. Neprohlašuj modelovou chybu za zjištěnou chybu třídy.

\begin{definiceBox}
Připomenutí: \textbf{negace výroku} má opačnou pravdivostní hodnotu
než původní výrok. Nestačí napsat jiné nepravdivé tvrzení;
negace musí pokrýt právě všechny případy, kdy původní výrok neplatí.
\end{definiceBox}

\begin{prikladReseniBox}[title={P1: společný rozbor dvou modelových chyb},breakable=false]
a) Žák negoval \uv{Každé celé číslo je kladné} větou
\uv{Každé celé číslo je záporné}. Oprav negaci.

b) Jiný žák přepsal $-1\le x<2$, $x\in\mathbb R$, jako $(-1;2\rangle$.
Oprav zápis a ověř oba krajní body.
\tcblower
a) Správně: \uv{Existuje celé číslo, které není kladné}, tedy
$\exists x\in\mathbb Z:\ x\le0$. Původní výrok je nepravdivý,
negace pravdivá, například díky $x=0$. Neplatnost obecného výroku
neznamená, že opačnou vlastnost mají všechny prvky.

b) $\langle-1;2)$. Číslo $-1$ vyhovuje, číslo $2$ nevyhovuje;
u každého kraje se proto musí změnit druh závorky.
\end{prikladReseniBox}

\blok{Navazující procvičování podle potřeby}
\poznamka{Z následujících dvou bloků zvol základ podle pozorovaných potíží.
Rychlejší žáci zpracují oba. Žák nejprve řeší samostatně, pak vysvětlí
spolužákovi první rozhodující krok. Spodní části slouží učiteli ke kontrole.}

\begin{zadaniReseniBox}[title={S1: znaménko, zlomek, interval},breakable=false]
\begin{enumerate}[label=\alph*)]
\item Vypočti $-18:(-3)-2\cdot[4-(-1)]$.
\item Porovnej $-\frac56$ a $-\frac45$ a ukaž společného jmenovatele.
\item Pro $x\in\mathbb R$ zapiš $-2<x\le0$ intervalem a nakresli osu.
Patří tam $-2$, $-\frac12$ a $0$?
\item Vypočti vzdálenost bodů $-3$ a $1$; je stejná jako rozdíl $-3-1$?
\end{enumerate}
\tcblower
a) $6-2\cdot5=-4$.
b) $-\frac{25}{30}<-\frac{24}{30}$, tedy $-\frac56<-\frac45$.
c) $(-2;0\rangle$: prázdný kraj $-2$, plný $0$; patří $-\frac12$ a $0$.
d) $|-3-1|=4$, ale rozdíl $-3-1=-4$; vzdálenost je nezáporná.
\end{zadaniReseniBox}

\begin{zadaniReseniBox}[title={S2: pravdivost a negace},breakable=false]
Nechť $A$ je výrok $0\in\mathbb N$ a $B$ výrok $-2\in\mathbb Z$.
Urči pravdivost $A\lor B$, $A\land B$, $A\Rightarrow B$, $B\Rightarrow A$,
$A\Leftrightarrow B$. Potom neguj $\exists x\in\mathbb Q:\ x>3$.
\tcblower
Při naší konvenci je $A$ nepravdivý, $B$ pravdivý.
Postupně: pravda, nepravda, pravda, nepravda, nepravda.
Implikace je nepravdivá jen při pravdivém předpokladu a nepravdivém závěru.
Negace: $\forall x\in\mathbb Q:\ x\le3$.
Původní existenční výrok je pravdivý například pro $x=4$, negace nepravdivá.
\end{zadaniReseniBox}

\blok{Rezerva a závěrečná zpětná vazba}
\begin{zadaniReseniBox}[title={R1: volitelná zásoba, 10–15 minut},breakable=false]
\begin{enumerate}[label=\alph*)]
\item Oprav výrok \uv{Každý nekonečný desetinný rozvoj je iracionální}
pomocí protipříkladu. Co k určení iracionality chybí?
\item Číslo $3{,}449$ zaokrouhli přímo na desetiny, potom přes setiny.
Vysvětli rozdíl.
\item Vyvrať $|a-b|=|a|-|b|$ vhodnou dvojicí čísel a vysvětli levý výraz.
\end{enumerate}
\tcblower
a) $\frac13=0{,}\overline3$ je racionální. Iracionální číslo má rozvoj
nekonečný a neperiodický, nikoli pouze nekonečný.
b) Přímo $3{,}4$; přes $3{,}45$ vyjde $3{,}5$. Zaokrouhlujeme z původního čísla.
c) Například $a=0,b=2$: vlevo $2$, vpravo $-2$.
Levá strana vyjadřuje vzdálenost bodů $a,b$.
\end{zadaniReseniBox}

\poznamka{Závěr: každý žák napíše jednu chybu, které nyní rozumí,
a jeden konkrétní správný příklad. Pro hodinu 11 si zaznamenej skutečné
typy obtíží bez osobních údajů. Nezaměňuj tuto zpětnou vazbu za další test.}

\blok{Domácí procvičování}
\begin{zadaniReseniBox}[title={D1: přibližně 10–15 minut},breakable=false]
\begin{enumerate}[label=\arabic*)]
\item Vypočti $-4-(-9)+2\cdot(-3)$ a vzdálenost bodů $-2$ a $5$.
\item Zkrať $\frac{-42}{63}$ a porovnej výsledek s $-\frac35$.
\item Zapiš intervalem $0\le x<4$, $x\in\mathbb R$, a zakresli ho.
\item Neguj $\forall x\in\mathbb Z:\ x\ne0$ a rozhodni o pravdivosti
původního výroku i negace.
\item Pro nepravdivý $A$ a pravdivý $B$ urči pravdivost
$A\Rightarrow B$ a $\neg(A\Rightarrow B)$. Vysvětli.
\end{enumerate}
\tcblower
1. $5-6=-1$; $|-2-5|=7$.
2. $-\frac23=-\frac{10}{15}<-\frac9{15}=-\frac35$.
3. $\langle0;4)$: plný kraj $0$, prázdný $4$.
4. $\exists x\in\mathbb Z:\ x=0$; původní nepravda, negace pravda.
5. Implikace je pravdivá, její negace nepravdivá:
nenastal případ pravdivého předpokladu a nepravdivého závěru.
\end{zadaniReseniBox}

\end{document}
